
\section{Conclusiones}

A lo largo del desarrollo del presente trabajo se logró caracterizar experimentalmente un regulador lineal LDO discreto, evaluando su comportamiento estático y dinámico mediante distintas metodologías de medición. La comparación entre los resultados experimentales y las simulaciones realizadas en LTspice permitió validar gran parte de los modelos teóricos empleados durante el diseño, observándose una buena correlación tanto en la regulación de tensión como en la respuesta frente a variaciones de carga y alimentación. Mientras que en la limitación de corriente las mediciones permitieron verificar una discrepancia con el diseño de 500 mA, lo cual es un error critico y no deseado.

A partir de la interpolación y análisis de las mediciones experimentales se verificó el correcto funcionamiento de la etapa de regulación dentro de la región esperada, incluyendo el comportamiento del \textit{dropout}, la regulación de línea, la regulación de carga y la eficiencia del sistema. Asimismo, los ensayos dinámicos mediante carga activa permitieron observar la respuesta transitoria del regulador frente a cambios abruptos de corriente, obteniéndose una dinámica subamortiguada coherente con la prevista en simulación, aunque con amplitudes superiores debido a efectos no ideales presentes en la implementación física.

Durante el desarrollo experimental también se identificó una discrepancia importante respecto de las simulaciones iniciales: en una primera instancia no se consideró adecuadamente el nodo correspondiente al limitador de corriente, lo cual introdujo una oscilación adicional  de aproximadamente 16 kHz no contemplada en el checkpoint previo. Este fenómeno evidenció la sensibilidad del sistema frente a polos y acoplamientos parásitos asociados a la etapa de limitación de corriente. A partir del análisis posterior se determinó que la incorporación de un capacitor de compensación de aproximadamente $100\ \text{pF}$ entre base y colector del transistor NPN del limitador permitiría estabilizar dicho nodo, reduciendo significativamente la tendencia oscilatoria observada.

En términos generales, el trabajo permitió no sólo validar experimentalmente los conceptos teóricos asociados al diseño de reguladores lineales LDO, sino también adquirir experiencia práctica en simulación analógica, compensación en frecuencia, caracterización experimental y análisis de estabilidad en sistemas electrónicos reales. La comparación entre teoría, simulación e implementación física resultó fundamental para comprender las limitaciones de los modelos ideales y la influencia de las no idealidades presentes en circuitos electrónicos discretos.